% !TeX encoding = UTF-8
\documentclass{../plantilla/hvtbook}
\HVTTitulo{Electrolizadores y seguridad}
\HVTSubtitulo{Comprende el principio de electrólisis, las variables que se miden y los criterios mínimos de seguridad antes de cualquier práctica supervisada con hidrógeno.}
\HVTNivel{Nivel técnico}\HVTSerie{Guía formativa}\HVTNumero{03}\HVTAccion{Empezar}
\HVTDescripcion{Guía formativa sobre electrolizadores, medición y seguridad.}
\HVTCanonical{https://hidrogenoverdeturquesa.com/curso-electrolizadores}
\HVTRegreso{Volver al catálogo}{/fundacion\#cursos}
\HVTPortada{../../images/Electrolizador.png}{Figura 1. Módulo de electrólisis para aprendizaje técnico supervisado.}
\begin{document}\HVTMakeTitle
\begin{HVTPage}{Capítulo 1}{Qué hace un electrolizador}{Objetivo de la lección}
\HVTLead{Al terminar, el estudiante explica el proceso general de electrólisis, identifica entradas y salidas del sistema, y reconoce los riesgos básicos asociados a electricidad, gas, presión y ventilación.}
\begin{HVTGrid}
\begin{HVTColumn}{Aprenderás}\begin{itemize}\item Qué variables entran al sistema: agua, electricidad y control. \item Qué salidas deben vigilarse: gases, calor y humedad. \item Por qué la ventilación y la detección son obligatorias. \item Qué significa práctica supervisada.\end{itemize}\end{HVTColumn}
\begin{HVTColumn}{Materiales de estudio}\begin{itemize}\item Imagen o ficha del electrolizador. \item Multímetro para clase demostrativa. \item Diagrama de flujo del sistema. \item Lista de riesgos y controles.\end{itemize}\end{HVTColumn}
\end{HVTGrid}
\HVTImagen{../../images/Electrolizador.png}{Figura 1-1. Puertos, conexiones y orientación física deben identificarse antes de cualquier operación.}
\end{HVTPage}
\begin{HVTPage}{Concepto}{Balance básico del sistema}{La idea central}
\HVTText{Un electrolizador transforma energía eléctrica y agua en gases. En una clase inicial no se busca maximizar producción, sino entender el sistema, medir de forma responsable y documentar condiciones de seguridad.}
\begin{HVTFlow}\HVTFlowItem{1}{Entrada}{Agua acondicionada y energía eléctrica.}\HVTFlowItem{2}{Proceso}{Separación electroquímica bajo control.}\HVTFlowItem{3}{Salida}{Gases, calor y datos de operación.}\end{HVTFlow}
\begin{HVTExample}{Ejemplo 3-1}{Lista mínima de revisión antes de una práctica}\begin{enumerate}\item Zona ventilada y sin fuentes de ignición. \item Conexiones eléctricas revisadas y secas. \item Recipientes, mangueras y salidas identificadas. \item Docente o responsable técnico presente. \item Plan de apagado inmediato.\end{enumerate}\end{HVTExample}
\HVTText{Comentario: el primer aprendizaje no es producir gas; es aprender a no operar a ciegas.}
\end{HVTPage}
\begin{HVTPage}{Práctica}{Mapa de variables}{Dibuja el sistema antes de operarlo}
\begin{HVTSteps}\item Marca entradas: agua, alimentación eléctrica y señal de control. \item Marca salidas: salida de gas, drenaje, calor y datos. \item Marca riesgos: electricidad, gas, presión, humedad y temperatura. \item Define mediciones: voltaje, corriente, temperatura y tiempo. \item Define control: cómo se apaga y quién autoriza la prueba.\end{HVTSteps}
\begin{HVTCode}{Registro de práctica}fecha:
responsable:
variable medida:
estado inicial:
observaciones:
acción de seguridad:\end{HVTCode}
\end{HVTPage}
\begin{HVTPage}{Actividad}{Evidencia de aprendizaje}{Ficha de seguridad del prototipo}
\begin{HVTTask}{Entrega}\item Dibuja el electrolizador como caja de proceso. \item Señala mínimo cinco riesgos y cinco controles. \item Explica por qué el hidrógeno requiere ventilación. \item Escribe el procedimiento de apagado en tres pasos.\end{HVTTask}
\end{HVTPage}
\end{document}
